% --- Άσκηση 36677 (36677.tex) ---
\Askhsh{36677}{4}
\begin{ylh}{2}
    \item 5.2. Αριθμητική πρόοδος
    \item 5.3. Γεωμετρική πρόοδος
\end{ylh}
Μια οικογένεια, προκειμένου να χρηματοδοτήσει τις σπουδές του παιδιού της, έχει να επιλέξει μεταξύ δυο προγραμμάτων που της προτείνονται:

Για το πρόγραμμα Α πρέπει να καταθέσει τον 1ο μήνα 1 ευρώ, το 2ο μήνα 2 ευρώ, τον 3ο μήνα 4 ευρώ και γενικά, κάθε μήνα που περνάει, πρέπει να καταθέτει ποσό διπλάσιο από αυτό που κατέθεσε τον προηγούμενο μήνα.

Για το πρόγραμμα Β πρέπει να καταθέσει τον 1ο μήνα 100 ευρώ, το 2ο μήνα 110 ευρώ, τον 3ο μήνα 120 ευρώ και γενικά, κάθε μήνα που περνάει πρέπει να καταθέτει ποσό κατά 10 ευρώ μεγαλύτερο από εκείνο που κατέθεσε τον προηγούμενο μήνα.

\begin{alist}
\item Να βρείτε
\begin{rlist}
\item το ποσό $a_\nu$ που πρέπει να κατατεθεί στο λογαριασμό τον ν\textsuperscript{ο} (νιοστό) μήνα σύμφωνα με το πρόγραμμα Α.

\monades{4}
\item το ποσό $\beta_\nu$ που πρέπει να κατατεθεί στο λογαριασμό τον ν\textsuperscript{ο} μήνα σύμφωνα με το πρόγραμμα Β.

\monades{4}
\item το ποσό $\mathrm{A}_\nu$ που θα υπάρχει στο λογαριασμό μετά από ν μήνες σύμφωνα με το πρόγραμμα Α.

\monades{5}
\item το ποσό $\mathrm{B}_\nu$ που θα υπάρχει στο λογαριασμό μετά από ν μήνες σύμφωνα με το πρόγραμμα Β.

\monades{5}
\end{rlist}
\item
\begin{rlist}
\item Τι ποσό θα υπάρχει στο λογαριασμό μετά τους πρώτους 6 μήνες, σύμφωνα με κάθε πρόγραμμα;
\monades{3}
\item Αν κάθε πρόγραμμα ολοκληρώνεται σε 12 μήνες, με ποιο από τα δύο προγράμματα το συνολικό ποσό που θα συγκεντρωθεί θα είναι μεγαλύτερο;
\monades{4}
\end{rlist}
\end{alist}
